\documentclass[12pt, a4paper]{article}

\usepackage[utf-8]{inputenc}
\usepackage[margin=1in]{geometry}
\usepackage{amsmath}
\usepackage{amssymb}
\usepackage{xcolor}
\usepackage{hyperref}
\usepackage{setspace}
\usepackage{enumitem}

% Set line spacing
\onehalfspacing

% Title and Author
\title{Spatial Disparity Between Scenic Potential and Road Accessibility in Northern Pakistan:\\
A Tehsil-Level Spatial Analysis}
\author{}
\date{March 2026}

\begin{document}

\maketitle

\section{Problem Statement}

Northern Pakistan (Gilgit-Baltistan and Upper Khyber Pakhtunkhwa) contains globally significant mountain landscapes but exhibits uneven tourism development. Infrastructure is concentrated along major corridors (e.g., Karakoram Highway), while potentially high-scenic regions remain remote and underdeveloped. Despite tourism's economic importance, no quantitative, spatially explicit framework systematically identifies disparities between scenic potential and accessibility at the Tehsil (sub-district) scale.

\section{Research Questions}

\begin{enumerate}
    \item Where are statistically significant clusters and spatial outliers of scenic potential located?
    \item Which Tehsils represent ``High Scenic--Low Accessibility'' priority regions for targeted development?
    \item Does accessibility meaningfully predict scenic potential variation across the region?
\end{enumerate}

\noindent\textbf{Significance:} Evidence-based answers support sustainable tourism planning, conservation prioritization, and balanced regional development policy.

\section{Methodology}

\subsection{Study Design}

Spatial analysis of 60--80 Tehsils using composite indices and local spatial clustering methods.

\subsection{Scenic Potential Index (SPI)}

Composite indicator combining z-score normalized components:
\begin{itemize}
    \item Terrain Ruggedness Index (SRTM 30m DEM; 40\% weight): Topographic complexity
    \item Forest cover (ESA WorldCover 2021; 25\% weight): Vegetation character
    \item Water bodies (ESA WorldCover; 20\% weight): Hydrological features
    \item Snow/Ice coverage (MODIS/Sentinel-2; 15\% weight): Glacier/peak aesthetics
\end{itemize}

\noindent\textbf{Formula:}
\[
\text{SPI} = 0.4 \times Z(\text{TRI}) + 0.25 \times Z(\text{Forest}) + 0.2 \times Z(\text{Water}) + 0.15 \times Z(\text{Snow})
\]

\subsection{Accessibility Index (AI)}

\[
\text{AI} = 0.5 \times Z(\text{Road Density}) - 0.5 \times Z(\text{Distance to Roads})
\]

Components derived from OpenStreetMap road network (primary/secondary roads only).

\subsection{Data Sources}

All free and open-access:
\begin{itemize}
    \item Administrative boundaries: GeoBoundaries ADM3 (Pakistan, 554 units)
    \item Elevation: SRTM DEM v3.0 (USGS)
    \item Land cover: ESA WorldCover 2021
    \item Road network: OpenStreetMap via OSMnx
    \item Snow extent: MODIS or Sentinel-2
\end{itemize}

\subsection{Spatial Analysis Techniques}

\begin{itemize}
    \item Queen contiguity weights matrix (row-standardized)
    \item Global Moran's $I$ test (spatial autocorrelation)
    \item Local Indicators of Spatial Association (LISA) for cluster identification (HH, LL, HL, LH classifications)
    \item Ordinary Least Squares regression: $\text{SPI} \sim \text{AI}$
    \item Conditional spatial lag model (if $n > 30$ and residual autocorrelation detected)
\end{itemize}

\subsection{Implementation}

Python with open-source libraries (geopandas, rasterio, libpysal, esda, scikit-learn). Reproducible Jupyter notebook workflow with persistent caching of external data.

\section{Objectives}

\begin{enumerate}
    \item Develop reproducible, open-source framework for measuring scenic potential and accessibility in mountain regions.
    \item Identify statistically significant high-scenic/low-accessibility Tehsils (realistic expectation: 5--15 units among 60--80).
    \item Quantify spatial association between accessibility and scenic potential; evaluate causal mechanisms.
    \item Provide evidence-based guidance for sustainable tourism and conservation planning.
\end{enumerate}

\section{Expected Outcomes}

\begin{enumerate}
    \item Tehsil-level SPI and AI datasets with geospatial visualization
    \item LISA cluster maps identifying scenic concentration areas and spatial outliers
    \item Regression model quantifying accessibility-scenery association
    \item Priority region identification for development feasibility assessment
    \item Fully reproducible Python workflow applicable to other mountain regions
    \item Technical report (5--10 pages) with findings and policy recommendations
\end{enumerate}

\section{Proof of Concept}

Pilot study (March 2026) on Gilgit region validated all methodological components. Results: strong local LISA clustering (2 HH clusters, 7 LL clusters), strong accessibility-scenery association ($R^2 \approx 0.85$, small-sample estimate), and identification of priority regions (e.g., Shigar). Workflow fully reproducible.

\end{document}
